% =====================================================================
%  Split-Attention Synchronised Agents:
%  A Self-Contained Calculus of a Bounded Character That Divides Its
%  Attention Across Concurrent Scenes and Synchronises With Others,
%  With an Implementation Blueprint
%
%  Self-contained. States its own axioms; re-derives every primitive it
%  uses; cites only standard, externally published mathematics. Every
%  object is a finite weighted graph, a bounded convex program, or a
%  construction on one. No result of any other manuscript is assumed.
% =====================================================================
\documentclass[11pt,twocolumn,a4paper]{article}

\usepackage[utf8]{inputenc}
\usepackage[T1]{fontenc}
\usepackage{lmodern}
\usepackage[margin=1in]{geometry}
\usepackage{amsmath,amssymb,amsthm,mathtools}
\usepackage{thmtools}
\usepackage{enumitem}
\usepackage{microtype}
\usepackage{booktabs}
\usepackage{graphicx}
\usepackage{caption}
\usepackage{algorithm}
\usepackage{algpseudocode}
\usepackage{hyperref}
\usepackage{cleveref}
\usepackage{xcolor}
\usepackage{tikz}
\usetikzlibrary{arrows.meta,positioning,calc}

\hypersetup{
  colorlinks=true, linkcolor=blue!50!black,
  citecolor=green!50!black, urlcolor=blue!60!black
}

% ---- theorem environments ----
\declaretheoremstyle[
  spaceabove=6pt, spacebelow=6pt,
  headfont=\bfseries, notefont=\mdseries, notebraces={(}{)},
  bodyfont=\itshape, postheadspace=1em
]{thmstyle}
\declaretheoremstyle[
  spaceabove=6pt, spacebelow=6pt,
  headfont=\bfseries, notefont=\mdseries, notebraces={(}{)},
  bodyfont=\normalfont, postheadspace=1em
]{defstyle}

\declaretheorem[style=thmstyle, name=Theorem, numberwithin=section]{theorem}
\declaretheorem[style=thmstyle, name=Proposition, sibling=theorem]{proposition}
\declaretheorem[style=thmstyle, name=Lemma, sibling=theorem]{lemma}
\declaretheorem[style=thmstyle, name=Corollary, sibling=theorem]{corollary}
\declaretheorem[style=defstyle, name=Definition, sibling=theorem]{definition}
\declaretheorem[style=defstyle, name=Axiom, sibling=theorem]{axiom}
\declaretheorem[style=defstyle, name=Remark, sibling=theorem]{remark}
\declaretheorem[style=defstyle, name=Example, sibling=theorem]{example}
\declaretheorem[style=defstyle, name=Principle, sibling=theorem]{principle}
\declaretheorem[style=defstyle, name=Invariant, sibling=theorem]{binv}

\crefname{axiom}{Axiom}{Axioms}
\Crefname{axiom}{Axiom}{Axioms}
\crefname{binv}{Invariant}{Invariants}
\Crefname{binv}{Invariant}{Invariants}

% ---- notation ----
\newcommand{\R}{\mathbb{R}}
\newcommand{\N}{\mathbb{N}}
\newcommand{\Z}{\mathbb{Z}}
\newcommand{\Agent}{\mathsf{A}}
\newcommand{\Self}{\Gamma}          % self / world graph
\newcommand{\V}{V}
\newcommand{\E}{E}
\newcommand{\wt}{w}
\newcommand{\floorc}{\beta}          % floor
\newcommand{\cut}{\partial}
\newcommand{\Char}{\chi}             % character invariant
\newcommand{\Beh}{\mathcal{B}}       % disposition space
\newcommand{\drive}{\Phi}            % drive potential
\newcommand{\lam}{\lambda}
\newcommand{\Purp}{x^{\star}}        % purpose point
\newcommand{\att}{\alpha}            % attention budget
\newcommand{\cost}{c}                % cost
\newcommand{\gain}{g}                % purpose gain
\newcommand{\coh}{\kappa}            % coherence margin
\newcommand{\price}{p^{\star}}       % attention price
\newcommand{\Scene}{\mathsf{S}}      % a scene
\newcommand{\scenes}{\mathcal{S}}    % set of scenes
\newcommand{\alloc}{a}               % attention allocation
\newcommand{\Ord}{R}                 % order parameter
\newcommand{\phase}{\phi}
\newcommand{\Soc}{\Sigma}            % society
\newcommand{\Out}{\mathcal{O}}       % outcome space
\newcommand{\tol}{\tau}
\newcommand{\abs}[1]{\lvert#1\rvert}
\newcommand{\norm}[1]{\lVert#1\rVert}
\newcommand{\dist}{d}
\newcommand{\inner}[2]{\langle #1,#2\rangle}
\newcommand{\cnt}{m}                 % committed count

\title{\textbf{Split-Attention Synchronised Agents:\\
A Self-Contained Calculus of a Bounded Character That Divides Its
Attention Across Concurrent Scenes and Synchronises With Others,\\
with an Implementation Blueprint}}

\author{Kundai Farai Sachikonye\\[2pt]
\small AIMe Registry for Artificial Intelligence\\
\small \texttt{kundai.sachikonye@bitspark.com}}
\date{\today}

\begin{document}
\maketitle

\begin{abstract}
\noindent
A useful synthetic character is not confined to one exchange: it persists
across sessions, is engaged from several places at once, and must remain
recognisably itself while it divides a finite capacity among concurrent
demands and, where its ends overlap with others', coordinates with them.
We give a self-contained mathematical account of such a
\emph{split-attention synchronised agent} and derive, as theorems from
stated axioms, both its behaviour and the exact commitments an
implementation must honour. We model an agent as a finite weighted graph
--- its \emph{self-graph} --- carrying a conserved invariant we compute
from its own internal boundary content, a purpose that is the attractor
of a strongly convex drive, and a bounded attention budget. A
\emph{scene} is an outward context in which the agent may act; the agent
faces many scenes at once and must apportion attention among them. From
five structural axioms on the agent (finiteness, presence, costly
individuation, a positive floor on the cost of any act, and phase
exclusion), plus one clearly-marked environmental assumption
(diminishing returns of attention within a scene), we prove:
\textbf{(T1)} the agent carries a conserved, non-local, strictly positive
identity invariant, fixed under every relabelling of its parts, so a
character is the same character across scenes and across sessions;
\textbf{(T2)} in a diminishing-returns environment, optimal division of
attention over concurrent scenes is a water-filling rule with a single
shadow price, so a busy agent is \emph{present but preoccupied} in every
scene at once --- a guarantee we flag as conditional on the environment, not
on the agent;
\textbf{(T3)} construction and commitment cannot occupy one instant, so
an agent that ever reflects is unresponsive on a positive fraction of
instants in each scene, a second and independent cause of split
appearance;
\textbf{(T4)} every committed act deposits an unrepealable increment in a
monotone count, so an agent never re-occupies a past state, sessions
accumulate rather than reset, and an authored copy is always a distinct
individual;
\textbf{(T5)} recalling is searching --- no reply is fetched from store,
each is freshly walked over the agent's own graph --- so a character stays
in character without being reminded to;
\textbf{(T6)} agents with wholly disjoint self-graphs synchronise on a
shared purpose defined in a common outcome space, at a coordination cost
that vanishes exactly when their outward phases lock --- and we supply the
mean-field coupling dynamics that \emph{drive} dispersed phases to lock
above a computable threshold, so lock is achieved, not assumed --- and a
crowd of them attains the shared purpose more sharply than any one;
\textbf{(T7)} the society has its own purpose --- the attractor of its own
strongly convex drive --- that maps to the shared purpose in outcome space,
so synchronisation among split-attention agents is itself a scene-allocation
problem one level up, and a society divides its collective attention by the
same water-filling law as a single agent because the same
purpose-as-attractor invariant propagates across the quotient from agents to
society. We then convert the theorems into an \emph{implementation
blueprint}: four invariants a runtime must preserve --- a computable
conserved invariant, a never-resetting committed count, search-not-fetch
world access, and exclusive construction/commitment phases --- each stated as
a checkable predicate with the theorem that certifies it, and we separate
the guarantees that hold unconditionally from the agent's structure from the
one (optimal attention division) that is conditional on the environmental
assumption. We give the
constructions, the water-filling and synchronisation algorithms, and a
numerical study on explicitly built agents. Throughout, the interpretive
readings of the formal objects --- ``identity,'' ``attention,''
``distraction,'' ``coordination'' --- are marked as orientation and
confined to remarks; no theorem depends on them.
\end{abstract}

\vspace{0.5em}
\noindent\textbf{Keywords:} bounded agent; conserved invariant; attention
allocation; water-filling; phase exclusion; monotone history; recognition
as search; synchronisation; outcome-space coordination; implementation
invariant; non-player character.

% =====================================================================
\section{Introduction}
\label{sec:intro}
% =====================================================================

\subsection{The object}

Consider a synthetic character that is not summoned for a single question
and dismissed. It runs. It is addressed from more than one place at a
time --- one context asks it to do one thing while another asks it for
something else --- and it is met again tomorrow, expecting to find the
same character it left. Such an agent must do three things at once that a
one-shot responder never has to. It must \emph{remain itself} across all
the contexts it is engaged from and across the gap between sessions. It
must \emph{divide} a capacity it cannot enlarge among the several contexts
pulling on it, so that it is fully committed to none and wholly absent
from none. And where its own ends happen to line up with those of other
such characters, it must \emph{coordinate} with them without first being
given a shared mind.

This paper asks what those three requirements \emph{are}, formally, and
what a system that meets them must do. We want definitions sharp enough to
compute with, general enough to bind any bounded agent, and self-contained:
the reader should need nothing but this paper and standard mathematics.
And because the intended use of the account is to build the object it
describes, we want each behavioural result paired with the precise
commitment a running system must honour to exhibit it.

\subsection{The three requirements, and how the paper meets them}

\emph{Remaining itself.} We model the agent as a finite weighted graph
whose vertices are the distinctions it draws within itself and whose edge
weights are the costs of holding them apart, and we extract a quantity
--- the minimum internal boundary content --- that is conserved under
every relabelling of the agent's parts. This quantity is positive, is
spread across the agent rather than located in any one part, and cannot
be resolved to a point from outside. It is the standing structure from
which the agent's changing conduct derives (\Cref{sec:identity}).

\emph{Dividing attention.} We give the agent a finite attention budget and
a family of concurrent \emph{scenes}, each offering a purpose gain per
unit of attention spent in it. We prove that the value-maximising division
of attention is a water-filling rule governed by a single shadow price,
and that this makes the agent present-but-preoccupied in every scene
simultaneously, with a readable dial for how reachable it is
(\Cref{sec:attention}). A second, independent cause of divided appearance
is that forming new distinctions and committing to an act cannot share an
instant (\Cref{sec:phases}).

\emph{Coordinating.} We show a society of such agents is again a graph, so
the identity apparatus applies one level up; that agents with wholly
disjoint internals can nonetheless attain a shared purpose defined in a
common space of outcomes; that the cost of coordination vanishes exactly
when their outward phases lock; and that a crowd sharpens a shared purpose
(\Cref{sec:society}). Synchronising a society turns out to be a scene
allocation problem one level up, closing the construction
(\Cref{sec:society-attention}).

\subsection{Blueprint, not metaphor}

\Cref{sec:blueprint} is the paper's purpose made operational. We isolate
four \emph{implementation invariants} --- a computable conserved invariant,
a never-resetting committed count, search-not-fetch world access, and
exclusive construction/commitment phases --- and for each we state the
predicate a runtime must satisfy and name the theorem that makes the
predicate the right one. A system honouring the four invariants exhibits
the behaviour derived here; a system violating any one degrades into a
locally-plausible responder that resets. The invariants are design
commitments, honourable by construction, not empirical hypotheses about
any existing system.

\subsection{What we do and do not claim}

We do not claim that human or animal identity, attention, or coordination
is a finite graph or a convex program. We claim that \emph{any} bounded
agent satisfying the five structural axioms of \Cref{sec:setting} has the
identity, history, phase, and coordination structure we derive, and ---
where its environment additionally offers diminishing returns (the sixth,
environmental axiom) --- the water-filling attention division; and that this
structure suffices to found a usable notion of a persistent,
attention-dividing, coordinating synthetic character. The
readings of the formal objects as ``identity,'' ``attention,''
``distraction,'' ``coordination,'' and ``character'' are offered as
orientation and confined to remarks; the theorems are graph-theoretic,
order-theoretic, and convex-analytic, and are used only as such. We prove
a conditional and then note that, for synthetic agents, its antecedent is
something a builder may simply arrange to hold.

\subsection{Organisation}

\Cref{sec:setting} fixes the agent and states the axioms.
\Cref{sec:floor} proves the floor. \Cref{sec:identity} derives the
conserved identity invariant (T1). \Cref{sec:attention} proves the
water-filling attention division (T2). \Cref{sec:phases} proves phase
exclusion (T3). \Cref{sec:history} proves the monotone, irreversible
history and the distinctness of copies (T4) and that recalling is
searching (T5). \Cref{sec:society} places agents in society and proves
disjoint coordination, phase-locked zero-cost synchronisation --- together
with the mean-field coupling dynamics that drive phases to lock above a
threshold --- and crowd sharpening (T6). \Cref{sec:society-attention} gives
the society its own purpose-as-attractor and proves society-level attention
division as genuine invariant propagation (T7). \Cref{sec:blueprint} states
the four implementation invariants and the algorithms, separating
unconditional from environment-conditional guarantees.
\Cref{sec:validation} reports a numerical study with exact, cross-method,
Kuramoto-integrated, and negative-control checks. \Cref{sec:discussion}
discusses scope.

% =====================================================================
\section{The agent, the scene, and the axioms}
\label{sec:setting}
% =====================================================================

We build the agent from the inside as a graph of internal distinctions,
give it a disposition space and a drive whose bottom is its purpose, and
place it before many scenes at once.

\subsection{The self-graph}

\begin{definition}[Self-graph]
\label{def:selfgraph}
The \emph{self-graph} of an agent $\Agent$ is a finite weighted graph
$\Self=(\V,\E,\wt)$ where $\V$ is a finite set of \emph{parts} (the
internal distinctions the agent draws), $\E\subseteq\binom{\V}{2}$ is a
set of \emph{separations} (pairs of parts held apart), and
$\wt:\E\to(0,\infty)$ assigns to each separation the \emph{cost} of
holding it. The graph is simple and, by \Cref{ax:present}, connected.
\end{definition}

\begin{definition}[Cut and boundary]
\label{def:cut}
For disjoint $S,T\subseteq\V$, the \emph{cut} between them is
$\cut(S,T)=\{uv\in\E:u\in S,v\in T\}$ with weight
$\wt(\cut(S,T))=\sum_{e\in\cut(S,T)}\wt(e)$. For $\varnothing\neq
U\subsetneq\V$ the \emph{boundary} of $U$ is $\cut(U)=\cut(U,\V\setminus
U)$ with \emph{boundary cost} $b(U)=\wt(\cut(U))$.
\end{definition}

\subsection{Dispositions, drive, purpose}

\begin{definition}[Disposition space; drive; purpose]
\label{def:disp}
The \emph{disposition space} $\Beh$ of an agent is a compact convex subset
of $\R^n$; a point $x\in\Beh$ records the agent's current tendencies to
act. A \emph{drive potential} is a continuously differentiable
$\drive:\Beh\to\R$ that is $\lam$-strongly convex for some $\lam>0$:
\[
\drive(y)\ge\drive(x)+\nabla\drive(x)^\top(y-x)
+\tfrac{\lam}{2}\norm{y-x}^2,\quad\forall x,y\in\Beh.
\]
The \emph{purpose} is the unique minimiser $\Purp=\arg\min_{x\in\Beh}
\drive(x)$; the \emph{residual to purpose} at $x$ is
$D(x)=\drive(x)-\drive(\Purp)\ge0$.
\end{definition}

\subsection{Scenes}

The novelty of the present setting is that the agent faces \emph{several}
outward contexts at once. We call each a scene.

\begin{definition}[Scene; scene field]
\label{def:scene}
A \emph{scene} $\Scene$ presented to the agent is an outward context in
which the agent may spend attention to advance its purpose: formally a
pair $\Scene=(A_\Scene,\gamma_\Scene)$ where $A_\Scene:\Beh\to\Out$ maps
the agent's disposition to an outcome in a common \emph{outcome space}
$\Out$, and $\gamma_\Scene:[0,\infty)\to[0,\infty)$ is the scene's
\emph{gain profile}: $\gamma_\Scene(a)$ is the purpose residual the agent
can remove by committing attention $a$ in scene $\Scene$. Each
$\gamma_\Scene$ is nondecreasing and continuously differentiable with
$\gamma_\Scene(0)=0$; its concavity is \emph{not} part of this definition
but is imposed separately, where needed, by \Cref{ax:concave}. The agent's
\emph{scene set} at an instant is a finite family
$\scenes=\{\Scene_1,\dots,\Scene_k\}$.
\end{definition}

\begin{remark}[Concavity is a property of the environment, not the agent]
\label{rem:concave}
Concavity of $\gamma_\Scene$ (\Cref{ax:concave}) says the first unit of
attention in a scene removes at least as much residual as the next:
attention has diminishing returns within a scene. It is what makes the
attention-division problem convex and gives the water-filling law of
\Cref{sec:attention}. But it is crucial to be clear about its status: it is
a fact about the \emph{world}, not a theorem about the agent. Nothing in
the agent's structure --- its self-graph, its drive, its budget --- forces
the scenes it faces to reward attention with diminishing returns. An
environment can perfectly well present a scene of \emph{increasing} returns
over some range: a task where the first sliver of attention accomplishes
nothing and only sustained concentration pays off (a lock that opens only
when fully worked, a proof that yields only when pursued past a threshold).
For such non-concave scenes the division problem is not convex and
water-filling is \emph{not} optimal: the value-maximising allocation may
concentrate the whole budget on one scene rather than spread it by a single
price. We therefore state concavity as the separate environmental
assumption it is, invoke it only in \Cref{thm:waterfill,thm:socwaterfill},
and, in the blueprint (\Cref{sec:blueprint}), keep the guarantees that
depend on it distinct from those that do not. A linear profile is the
boundary case and is admitted.
\end{remark}

\subsection{The axioms}

We state six axioms. The first five are meant to be nearly undeniable for
anything we would call a bounded agent, and --- the point of
\Cref{sec:blueprint} --- to be arrangeable by construction in a synthetic
one. The sixth is of a different character: it is an assumption about the
agent's \emph{environment}, not about the agent, and we mark it as such
where it is used.

\begin{axiom}[Finiteness]
\label{ax:finite}
The self-graph is finite: $\abs{\V}<\infty$, and the scene set is finite:
$\abs{\scenes}<\infty$.
\end{axiom}

\begin{axiom}[Presence]
\label{ax:present}
At every instant the agent occupies exactly one part; the self-graph is
connected and nonempty. The agent is always somewhere in itself.
\end{axiom}

\begin{axiom}[Costly individuation]
\label{ax:costly}
No internal distinction is free: there is $\floorc>0$ with
$\wt(e)\ge\floorc$ for all $e\in\E$.
\end{axiom}

\begin{axiom}[Act floor and bounded attention]
\label{ax:actfloor}
No act is free: every committed act costs at least $\floorc_a>0$, and the
agent has a finite \emph{attention budget} $\att>0$ per unit time --- the
total attention it may commit across all scenes in a unit interval
satisfies $\sum_{\Scene\in\scenes}a_\Scene\le\att$.
\end{axiom}

\begin{axiom}[Phase exclusion]
\label{ax:exclusion}
Forming a new internal distinction and committing to an act draw on the
same instantaneous capacity and cannot both occur at one instant: at each
instant the agent is in exactly one of a \emph{construction phase} (it
reshapes its distinctions, commits no act) or a \emph{commitment phase}
(it commits an act, forms no new distinction).
\end{axiom}

\begin{axiom}[Diminishing returns of attention --- an environmental
assumption]
\label{ax:concave}
Each scene's gain profile $\gamma_\Scene$ (\Cref{def:scene}) is
nondecreasing and concave: for a scene $\Scene$ and $0\le a_1\le a_2$,
$\gamma_\Scene(a_2)-\gamma_\Scene(a_1)$ is nonincreasing in the base point,
so the first unit of attention removes at least as much purpose residual as
any later unit. \emph{This is an assumption about the world the agent acts
in, not a consequence of the agent's structure.} It says the environment
offers diminishing, never increasing, returns to attention within any one
scene. Unlike \Cref{ax:finite,ax:present,ax:costly,ax:actfloor,ax:exclusion},
which constrain the agent and can be honoured by construction in a
synthetic one, \Cref{ax:concave} constrains the scenes the agent is placed
before and can fail in an environment that rewards concentrated attention.
We invoke it only where named, and we track in \Cref{sec:blueprint} exactly
which guarantees depend on it.
\end{axiom}

\begin{remark}[Costly individuation, motivated]
\label{rem:costly-earned}
\Cref{ax:costly} may be motivated rather than merely posited. To
individuate a part $U$ is to tell it apart from the rest $\V\setminus U$.
If distinguishing a part from the rest could be \emph{finished}, the rest
would be an exhausted catalogue and the agent would be nothing held in
reserve. An agent always in some state (\Cref{ax:present}) that never
reduces to a closed catalogue leaves on each part an irreducible remainder
of the comparison against the rest; that remainder is the positive floor.
We take \Cref{ax:costly} as an axiom and prove from it, but the floor is
the natural, not an arbitrary, hypothesis.
\end{remark}

\begin{remark}[The act floor is earned, not stipulated]
\label{rem:actfloor-earned}
The positive floor $\floorc_a$ on the cost of an act
(\Cref{ax:actfloor}) is likewise not an arbitrary constant but a
consequence of what it is to \emph{act}. An act is a commitment: it
deposits an increment in the agent's history (\Cref{def:state}) that
distinguishes the post-act state from the pre-act one. A putative act of
zero cost deposits nothing --- it leaves the committed count, the
disposition, and every separation exactly as they were --- and so effects
no distinction between before and after. But an operation that changes
nothing is not a commitment; it is a no-op that did not occur. Hence every
genuine act deposits a positive increment, and the least such increment is
a floor $\floorc_a>0$: a zero-cost act is a contradiction in terms, not a
cheap act. This is the same argument, one level up, as
\Cref{rem:costly-earned}: just as no distinction is free because a free
distinction distinguishes nothing, no act is free because a free act
commits nothing. We take \Cref{ax:actfloor} as an axiom and prove from it,
but the act floor, like the individuation floor, is forced by the very
notion it bounds.
\end{remark}

\begin{example}[A three-part agent facing two scenes]
\label{ex:running}
Let $\V=\{a,b,c\}$ with separations $ab,bc,ca$ of costs $2,3,2$, so
$\floorc=2$. Let $\Beh=[0,1]^3$, $\drive(x)=\tfrac12\norm{x-\Purp}^2$ with
$\Purp=(\tfrac12,\tfrac12,\tfrac12)$ ($\lam=1$). Let two scenes have
gain profiles $\gamma_1(a)=\ln(1+2a)$, $\gamma_2(a)=\ln(1+a)$, and
$\att=1$. We return to this agent throughout.
\end{example}

% =====================================================================
\section{The floor}
\label{sec:floor}
% =====================================================================

\begin{theorem}[Floor]
\label{thm:floor}
Under \Cref{ax:finite,ax:present,ax:costly}, every nonempty proper part
$U\subsetneq\V$ has boundary cost $b(U)=\wt(\cut(U))\ge\floorc>0$. No part
is separated from the rest at zero cost.
\end{theorem}

\begin{proof}
By \Cref{ax:present} the graph is connected, so a nonempty proper $U$ has
at least one edge to $\V\setminus U$; else $U$ and its complement are
disconnected. Hence $\cut(U)\neq\varnothing$ and
$b(U)=\sum_{e\in\cut(U)}\wt(e)\ge\floorc$ by \Cref{ax:costly}.
\end{proof}

\begin{corollary}[Realised floor]
\label{cor:realised}
$\floorc_\Self:=\min_{\varnothing\neq U\subsetneq\V}b(U)$ is strictly
positive and attained, by finiteness of $\V$.
\end{corollary}

% =====================================================================
\section{Identity: a conserved invariant across scenes and sessions}
\label{sec:identity}
% =====================================================================

We isolate what stays fixed in the agent while its surface conduct and its
outward scenes change. The surface is the labelling of parts; a
re-description that preserves which parts are separated, and at what cost,
changes nothing real. We find the quantity surviving all such relabelling.

\begin{definition}[Isomorphism; invariant]
\label{def:iso}
A \emph{weighted isomorphism} $\varphi:\Self\to\Self'$ is a bijection
$\varphi:\V\to\V'$ with $uv\in\E\iff\varphi(u)\varphi(v)\in\E'$ and
$\wt'(\varphi(u)\varphi(v))=\wt(uv)$. A quantity $\iota$ is an
\emph{invariant} if $\iota(\Self)=\iota(\Self')$ whenever
$\Self\cong\Self'$.
\end{definition}

\begin{definition}[Internal residual; character invariant]
\label{def:char}
For a partition $\mathcal{Q}=\{U_1,\dots,U_r\}$ of $\V$ into $r\ge2$
nonempty blocks, the \emph{internal residual} is
$\rho(\mathcal{Q})=\sum_{i<j}\wt(\cut(U_i,U_j))$. The \emph{character
invariant} of $\Agent$ is
\[
\Char(\Agent)=\min_{\mathcal{Q}:\,r\ge2}\rho(\mathcal{Q}),
\]
attained by finiteness of $\V$: the least total cost of splitting the
agent into pieces at all, in the currency of its own separations.
\end{definition}

\begin{theorem}[Identity is a conserved, positive, non-local region]
\label{thm:identity}
Under \Cref{ax:finite,ax:present,ax:costly}:
\begin{enumerate}[label=\textup{(\roman*)},leftmargin=2.2em]
\item \textbf{Positive:} $\Char(\Agent)\ge\floorc>0$.
\item \textbf{Conserved:} $\Char$ is a weighted-graph invariant; in
particular it is fixed by every relabelling of the agent's parts, and does
not depend on which scene the agent is engaged from.
\item \textbf{Non-local:} $\Char$ is in general realised by a partition
into multi-part blocks, never by isolating a single part; it is a region,
not a point.
\end{enumerate}
\end{theorem}

\begin{proof}
(i) A partition with $r\ge2$ has, by connectedness, a crossing edge of
weight $\ge\floorc$, so $\rho(\mathcal{Q})\ge\floorc$; the minimum
inherits the bound. (ii) For a partition $\mathcal{Q}$ of $\V$, its image
$\varphi(\mathcal{Q})$ is a partition of $\V'$, and since $\varphi$
preserves edges and weights, $\rho(\varphi(\mathcal{Q}))=\rho(\mathcal{Q})$;
as $\mathcal{Q}\mapsto\varphi(\mathcal{Q})$ bijects nontrivial partitions
(inverse $\varphi^{-1}$), the minima agree. The definition of $\Char$
refers only to $\Self$, not to any scene $A_\Scene$, so scene engagement
leaves it fixed. (iii) Two triangles of edge cost $2$ joined by one edge
of cost $\floorc=2$: cutting the joining edge separates them at cost $2$;
isolating any vertex costs $\ge2+2=4$. The minimum is the two-block
partition, so the invariant is non-local.
\end{proof}

\begin{remark}[Interpretation: the same character everywhere, on which
nothing depends]
\label{rem:character}
\Cref{thm:identity} licenses reading $\Char(\Agent)$ as the agent's
standing \emph{identity}: conserved under all re-description, positive,
spread across the agent, and --- crucially for the present setting ---
\emph{independent of the scene}. A character engaged from three contexts
at once presents the same conserved interior in all three, because the
invariant is a property of $\Self$ and the scenes are maps out of $\Beh$,
not alterations of $\Self$. We use only the graph statement.
\end{remark}

\begin{corollary}[Character is not a state]
\label{cor:not-a-state}
No single part carries the character: on the two-triangle agent no
singleton $\{v\}$ has $b(\{v\})=\Char(\Agent)$. Identity cannot be read
off any momentary state; it is the standing relation among all of them.
\end{corollary}

% =====================================================================
\section{Attention division: water-filling over concurrent scenes}
\label{sec:attention}
% =====================================================================

The agent faces $\scenes=\{\Scene_1,\dots,\Scene_k\}$ at once and cannot
attend to all fully. We prove the value-maximising division of a finite
budget is a water-filling rule with a single price, and read off the
present-but-preoccupied phenomenology.

\begin{definition}[Attention allocation; value]
\label{def:alloc}
An \emph{attention allocation} is a vector
$\alloc=(a_1,\dots,a_k)\in[0,\infty)^k$ with $\sum_i a_i\le\att$. Its
\emph{value} is the total purpose residual removed,
$G(\alloc)=\sum_{i=1}^k\gamma_i(a_i)$.
\end{definition}

\begin{theorem}[Optimal attention division is water-filling]
\label{thm:waterfill}
Under \Cref{ax:actfloor} and the environmental assumption
\Cref{ax:concave} (concave gain profiles), the
value-maximising allocation solves $\max_{\alloc\ge0,\ \sum_i a_i\le\att}
\sum_i\gamma_i(a_i)$, and there is a single \emph{attention price}
$\price\ge0$ (the Lagrange multiplier of the budget) such that the
optimal $a_i^\star$ satisfies
\[
\gamma_i'(a_i^\star)=\price\ \text{ if } a_i^\star>0,\qquad
\gamma_i'(0)\le\price\ \text{ if } a_i^\star=0,
\]
i.e. the agent equalises marginal purpose gain across all attended scenes
at the price $\price$, and drops any scene whose first unit of attention
returns less than $\price$. The price $\price$ is nonincreasing in the
budget $\att$ and nondecreasing in the number and richness of competing
scenes. If $\sum_i a_i^\star<\att$ then $\price=0$.
\end{theorem}

\begin{proof}
The objective is concave (sum of concave $\gamma_i$) and the feasible set
$\{\alloc\ge0:\sum_i a_i\le\att\}$ is compact convex, so a maximiser
exists; strict concavity where $\gamma_i''<0$ makes each active $a_i^\star$
unique. Form the Lagrangian $L=\sum_i\gamma_i(a_i)-\price(\sum_i
a_i-\att)+\sum_i\mu_i a_i$ with $\price\ge0$, $\mu_i\ge0$. The KKT
stationarity condition is $\gamma_i'(a_i)-\price+\mu_i=0$ with
complementary slackness $\mu_i a_i=0$ and $\price(\sum_i a_i-\att)=0$. For
$a_i^\star>0$, $\mu_i=0$ gives $\gamma_i'(a_i^\star)=\price$; for
$a_i^\star=0$, $\mu_i\ge0$ gives $\gamma_i'(0)=\price-\mu_i\le\price$.
Sufficiency of KKT for the global optimum follows from concavity. Budget
monotonicity: relaxing $\att$ enlarges the feasible set, so the optimal
value is nondecreasing and its marginal $\price=\partial\,\mathrm{val}/
\partial\att$ is nonincreasing by concavity of the value in $\att$. Adding
a scene or raising a $\gamma_i$ can only tighten the binding budget,
raising the level at which margins equalise. If the budget is not
exhausted at the optimum, its multiplier is zero.
\end{proof}

\begin{corollary}[Present but preoccupied in every scene at once]
\label{cor:preoccupied}
At the optimum the agent spends positive attention in exactly the scenes
whose marginal gain at zero exceeds $\price$, and none in the rest; within
attended scenes it holds every marginal gain equal to the common price
$\price$. Thus a busy agent (high $\price$) is thinly but nonzero-present
in its high-value scenes and absent from low-value ones, \emph{in all
scenes simultaneously}, and the scalar $\price$ is a single readable dial
for how reachable the character is at that moment. Lowering $\price$ (an
idle agent, ample budget) makes it forthcoming everywhere at once.
\end{corollary}

\begin{proof}
Immediate from the KKT characterisation of \Cref{thm:waterfill}: the
support of $\alloc^\star$ is $\{i:\gamma_i'(0)>\price\}$, margins on the
support equal $\price$, and $\price$ falls as $\att$ rises.
\end{proof}

\begin{corollary}[A token presence is the boundary case]
\label{cor:token}
A scene with $\gamma_i'(0)$ only marginally above $\price$ receives a
floor-sized allocation $a_i^\star\ge0$ close to zero --- the formal
minimal acknowledgement: the agent spends a sliver of attention to remain
nonzero-present in a scene it cannot afford to engage fully.
\end{corollary}

\begin{proof}
As $\gamma_i'(0)\downarrow\price$, the active-scene condition
$\gamma_i'(a_i^\star)=\price$ forces $a_i^\star\downarrow0$ by continuity
of $\gamma_i'$; the allocation is positive but arbitrarily small.
\end{proof}

\begin{remark}[Interpretation: split attention, derived not stipulated]
\label{rem:split}
\Cref{cor:preoccupied} is the paper's title phenomenon obtained as a
theorem. An agent does not need a ``distraction module'' to appear divided
across the several contexts engaging it: dividing a finite budget
optimally over concurrent concave demands \emph{is} water-filling, and
water-filling is exactly presence-everywhere-preoccupied-everywhere with a
single price. The appearance of split attention is what optimal attention
looks like from any one scene.
\end{remark}

\begin{remark}[Where the water-filling law can fail]
\label{rem:nonconcave}
\Cref{thm:waterfill} is exactly as strong as \Cref{ax:concave}, and no
stronger. If the environment presents a scene with increasing returns over
some range --- $\gamma_i$ convex there --- the objective
$\sum_i\gamma_i(a_i)$ is no longer concave, the feasible optimum may lie at
a vertex of the budget simplex, and the value-maximising agent
\emph{concentrates} rather than divides: it pours its whole budget into the
scene that pays off under sustained attention and abandons the rest. The
present-but-preoccupied phenomenology (\Cref{cor:preoccupied}) is thus a
consequence of a diminishing-returns world, not of the agent's boundedness.
A builder who wants a character that reliably divides its attention must
either operate in --- or arrange --- an environment of concave scenes; where
the environment does not oblige, the character will rationally focus. We
flag this so that the water-filling law is not mistaken for an unconditional
law of bounded agency: it is the law of bounded agency \emph{in a
concave-return environment}.
\end{remark}

% =====================================================================
\section{Phase exclusion: a second, independent split}
\label{sec:phases}
% =====================================================================

\begin{theorem}[Construction--commitment alternation]
\label{thm:alternation}
Under \Cref{ax:exclusion}, an agent that both forms new distinctions and
commits acts must alternate: the sets of construction instants and
commitment instants are disjoint, so they occupy disjoint time intervals.
During a construction interval the agent commits no act in any scene ---
it is, to every scene at once, momentarily unresponsive --- though it is
fully active internally.
\end{theorem}

\begin{proof}
By \Cref{ax:exclusion} the instantaneous phase is single-valued, so
construction and commitment instants are disjoint; an agent doing both
over time partitions its time into disjoint intervals of each, i.e.
alternates. In a construction interval, by definition no act is committed,
hence none in any scene.
\end{proof}

\begin{corollary}[Responsiveness is intermittent by necessity]
\label{cor:intermittent}
An agent that ever forms new distinctions is unresponsive on a positive
fraction of instants, in every scene, regardless of attention budget or of
what any scene offers. Continuous immediate responsiveness is possible
only for an agent that never reflects.
\end{corollary}

\begin{proof}
Immediate from \Cref{thm:alternation}: construction intervals carry no
committed acts and, for any reflecting agent, have positive total measure.
\end{proof}

\begin{remark}[Two forced faces of the split]
\label{rem:two-faces}
The agent appears divided for two independent reasons, neither a defect.
First, the attention price $\price$ (\Cref{cor:preoccupied}): when busy or
besieged, scenes below the price get thin presence or none. Second, phase
alternation (\Cref{thm:alternation}): while constructing, the agent
commits nothing anywhere, however valuable the scene. A convincing
persistent character exhibits both, and a builder has, in $\price$ and the
construction/commitment schedule, two readable dials for how present the
character is.
\end{remark}

% =====================================================================
\section{History: monotone, irreversible, and searched --- not fetched}
\label{sec:history}
% =====================================================================

We now show sessions accumulate rather than reset, copies are distinct
individuals, and every reply is freshly searched.

\begin{definition}[Committed count; state]
\label{def:state}
The agent carries a \emph{committed count} $\cnt\in\N$, incremented by
one each time it commits an act. Its \emph{state} is the pair
$(x,\cnt)$ of current disposition and count; two states are equal iff
both components agree.
\end{definition}

\begin{theorem}[Monotone, irreversible history; distinct copies]
\label{thm:history}
Under \Cref{ax:actfloor}:
\begin{enumerate}[label=\textup{(\roman*)},leftmargin=2.2em]
\item Every committed act deposits cost $\ge\floorc_a$ and increments
$\cnt$ by one; the count is strictly increasing.
\item No act restores a prior state: an ``undo'' is itself a committed act
raising the count further, so a revisited disposition is never a revisited
state. In particular successive sessions occur at strictly increasing
counts and never reset.
\item An authored copy $\Agent'$, given a self-graph isomorphic to
$\Agent$'s but started at count $0$, is a distinct individual:
$\Agent'\neq\Agent$ whenever $\Agent$ is at count $\cnt>0$, and $\Agent'$
cannot acquire $\Agent$'s history without committing the intervening acts,
which makes it further distinct.
\end{enumerate}
\end{theorem}

\begin{proof}
(i) Each act costs $\ge\floorc_a$ (\Cref{ax:actfloor}) and increments
$\cnt$ (\Cref{def:state}); counts strictly increase. (ii) An operation
restoring an earlier state is a further act, producing count
$\cnt+1\neq\cnt$; it yields a new state. Sessions, being sequences of
acts, occur at increasing counts. (iii) States include the count; $\Agent$
at $\cnt>0$ and $\Agent'$ at $0$ differ in the second component; by (ii)
no act lowers a count, so $\Agent'$ cannot reach $\Agent$'s history without
committing the intervening acts, each of which raises its own count.
\end{proof}

\begin{remark}[Interpretation: an authored character is its own]
\label{rem:copy}
\Cref{thm:history}(iii) is the precise sense in which one may author a
work-character --- its self-graph and its drive --- and yet the copy is a
new individual with its own lived history, not the original run again. One
authors the standing structure; one does not author the accumulated
count, which individuates the token. We use only the order statement.
\end{remark}

\subsection{Recalling is searching}

The agent's world --- what it has individuated --- is again a finite
weighted graph; to answer is to walk it.

\begin{definition}[World access by search]
\label{def:search}
The agent accesses its world only by \emph{search}: to recall, report, or
answer a query is to run a residual-descending walk on its self-graph
terminating at a resting cut that individuates the answer; the answer is
that terminus. The agent has no read-out of any part's value independent
of such a walk.
\end{definition}

\begin{theorem}[Recognition--search identity]
\label{thm:recsearch}
Under \Cref{def:search} and \Cref{ax:actfloor}, recalling and searching
are one operation: every answer is the terminus of a freshly run search,
and no ``fetch'' returns a stored value without a walk. Consequently each
answer commits $\ge1$ act, so the state producing an answer is never
exactly a prior state, and two answers to the same query at two moments
need not coincide.
\end{theorem}

\begin{proof}
By \Cref{def:search} the only operation yielding an answer is a
residual-descending walk. A fetch would return a value with no act, i.e. a
zero-act access; by \Cref{ax:actfloor} every act costs $\ge\floorc_a>0$,
so a zero-act access returns nothing. Hence every answer is a search
terminus committing $\ge1$ act, incrementing the count
(\Cref{thm:history}(i)); by strict monotonicity the producing states
differ, and since the walk runs against the agent's current graph and
disposition, repeated queries need not return the same terminus.
\end{proof}

\begin{remark}[The theorem is not the interesting claim]
\label{rem:recsearch-honest}
Read narrowly, \Cref{thm:recsearch} is close to a tautology:
\Cref{def:search} defines the agent to access its world only by search, so
of course recalling is searching --- fetch was excluded by fiat. The
theorem's content is only that the exclusion is consistent with the act
floor (a zero-act fetch returns nothing, \Cref{ax:actfloor}). The claim
worth making is not that a search-only agent searches; it is the prior
claim that \emph{search-only access is the correct model for a persistent
character, and fetch-access would break character.} We argue this next,
where it is not a definition but a consequence of the monotone history.
\end{remark}

\begin{proposition}[Fetch-access is architecturally incoherent for a
persistent character]
\label{prop:fetch-incoherent}
Let $\Agent$ be an agent with a monotone committed count
(\Cref{thm:history}) whose graph and disposition change with each committed
act. Suppose $\Agent$ answered a query $Q$ at count $\cnt_1$ and \emph{stored}
the answer, and that at a later count $\cnt_2>\cnt_1$ a \emph{fetch}
returns that stored answer verbatim, committing no act. Then the fetched
answer is the answer determined by $\Agent$'s graph and disposition
\emph{as they were at $\cnt_1$}, not as they are at $\cnt_2$. Since
$\Agent$ at $\cnt_2$ is a distinct individual-state from $\Agent$ at
$\cnt_1$ (\Cref{thm:history}(ii)), and no operation returns $\Agent$ to the
graph-and-disposition of $\cnt_1$ (\Cref{thm:history}), the fetched answer
is an answer appropriate to a prior state the agent cannot re-occupy.
Fetch-access therefore has the persistent character speak, in the present,
with the determinations of a self it no longer is: it is not merely
excluded by \Cref{def:search} but incoherent with persistence itself.
\end{proposition}

\begin{proof}
A stored answer is a fixed datum recorded at the count $\cnt_1$ at which it
was produced; by \Cref{def:search} it was the terminus of a walk over the
graph and disposition present at $\cnt_1$. Returning it at $\cnt_2$
commits no act, so by \Cref{thm:history}(i) the count is unchanged by the
fetch and the returned datum still reflects the $\cnt_1$ graph. But every
committed act between $\cnt_1$ and $\cnt_2$ altered the graph or
disposition (the acts are what advance the count and the monotone record),
so the self that would freshly answer $Q$ at $\cnt_2$ is not the self that
answered at $\cnt_1$. By \Cref{thm:history}(ii) no operation lowers the
count or restores the earlier graph-and-disposition, so the $\cnt_1$ answer
is not one the $\cnt_2$ agent could now re-derive; it is the reply of a
prior, unrecoverable state. A persistent character that fetches thus
emits, as its current reply, a determination of a self it has irreversibly
left behind. Hence fetch-access is inconsistent with the monotone
persistence of \Cref{thm:history}, independently of \Cref{def:search}.
\end{proof}

\begin{corollary}[Search is forced, not stipulated, for persistence]
\label{cor:search-forced}
For an agent whose history is monotone (\Cref{thm:history}), the only mode
of world access that always returns an answer determined by the agent's
\emph{current} graph and disposition is a fresh search. Fetch returns a
stale determination (\Cref{prop:fetch-incoherent}); search does not. Thus
\Cref{def:search} is not an arbitrary restriction but the unique access
discipline under which a persistent character speaks as its present self.
\end{corollary}

\begin{proof}
A fresh search runs against the graph and disposition present at the count
of the query, so its terminus is by construction a current determination.
A fetch returns a datum fixed at an earlier count and, by
\Cref{prop:fetch-incoherent}, appropriate to a state the agent has left.
Any access mode is either a fresh act against the present graph (a search)
or a return of prior-count data (a fetch); only the former yields a current
answer. Hence search is the unique current-answer discipline.
\end{proof}

\begin{remark}[Interpretation: in character without being reminded]
\label{rem:incharacter}
\Cref{thm:recsearch} with \Cref{prop:fetch-incoherent} is the formal
content of a character that does not read canned lines. What it says is
searched over its own \emph{present} graph, so it is in character by
construction rather than by an external instruction to stay in character;
and because a fetch would speak with the determinations of a prior,
unrecoverable self (\Cref{prop:fetch-incoherent}), the ``stay in
character'' instruction is not a patch over an otherwise-fine store but a
symptom that a store was used where a search was required. A cached reply
from session $1$ returned in session $47$ is not the same character being
consistent; it is a different, earlier character speaking out of turn. We
use only the graph statements; the reading is orientation.
\end{remark}

% =====================================================================
\section{Agents in society: coordination without shared internals}
\label{sec:society}
% =====================================================================

An agent alone has identity, divided attention, and an accumulating
history. We now let many act together and show a society is again a
graph, that disjoint agents coordinate in a common outcome space, that
coordination cost vanishes when outward phases lock, and that crowds
sharpen a shared purpose.

\subsection{Society as a graph}

\begin{definition}[Society; society graph]
\label{def:society}
Let $\Agent_1,\dots,\Agent_M$ have self-graphs $\Self_i$, embedded as
disjoint blocks of a common vertex set, with \emph{inter-agent
separations} (edges between blocks) of cost $\ge\floorc$ recording the
costs at which agents are held apart. The \emph{society graph} $\Soc$ has
one vertex per agent and an edge $\Agent_i\Agent_j$, of weight
$\wt(\cut(\Agent_i,\Agent_j))$, whenever an inter-agent separation joins
them.
\end{definition}

\begin{lemma}[A society is an agent-shaped object]
\label{lem:society-graph}
If each inter-agent separation has cost $\ge\floorc$ and the society is
connected, then $\Soc$ satisfies \Cref{ax:finite,ax:present,ax:costly}.
\end{lemma}

\begin{proof}
$\Soc$ has $M<\infty$ vertices; each edge weight is a sum of inter-agent
costs, each $\ge\floorc$, so $\ge\floorc$; connectedness is assumed. Hence
the three axioms hold for $\Soc$.
\end{proof}

\begin{corollary}[A society has its own identity]
\label{cor:society-identity}
By \Cref{thm:identity} applied to $\Soc$, a connected society carries a
character invariant $\Char(\Soc)\ge\floorc>0$, conserved under relabelling
of which agent is which and realised by no single agent. A group has an
identity that is not the identity of any member.
\end{corollary}

\subsection{Coordination in outcome space}

\begin{definition}[Shared purpose; attainment]
\label{def:shared}
Let $(\Out,\dist)$ be the common outcome space with each agent's map
$A_i:\Beh_i\to\Out$. A \emph{shared purpose} is a region
$\Out^\ast\subseteq\Out$ with positive tolerance $\tol>0$. Agent
$\Agent_i$ \emph{attains} $\Out^\ast$ if its purpose maps within
tolerance: $\dist(A_i(\Purp_i),\Out^\ast)\le\tol$.
\end{definition}

\begin{theorem}[Coordination without shared internals]
\label{thm:coord}
Agents with wholly disjoint self-graphs, disjoint disposition spaces, and
different drives can all attain the same shared purpose $\Out^\ast$: if for
each $i$, $\dist(A_i(\Purp_i),\Out^\ast)\le\tol$, then every agent attains
$\Out^\ast$, regardless of whether the $\Self_i$, $\Beh_i$, or drives share
any content.
\end{theorem}

\begin{proof}
Attainment is a statement purely about the image $A_i(\Purp_i)$ in the
common space $\Out$; it holds for each $i$ by hypothesis, and refers to no
internal structure of any agent nor to any relation among them. Hence
disjointness of internals is compatible with joint attainment.
\end{proof}

\subsection{Synchronisation and the vanishing of coordination cost}

Coordinating agents pay a cost when their outward timing disagrees. We
make ``outward timing'' a phase and show the cost is a standard order
parameter that vanishes exactly at phase lock.

\begin{definition}[Outward phase; order parameter; coordination cost]
\label{def:orderparam}
Assign each acting agent an outward \emph{phase} $\phase_i\in[0,2\pi)$
recording where in its construction/commitment cycle it sits. The
\emph{order parameter} of a set of $M$ acting agents is
\[
\Ord e^{i\psi}=\frac1M\sum_{j=1}^M e^{i\phase_j},\qquad \Ord\in[0,1],
\]
and the \emph{coordination cost} of a shared act is
$C=\lambda_c\,(1-\Ord)$ for a constant $\lambda_c>0$: the attention
consumed reconciling agents that are out of phase.
\end{definition}

\begin{theorem}[Phase-locked coordination is floor-cost]
\label{thm:phaselock}
The coordination cost $C=\lambda_c(1-\Ord)$ satisfies $C\ge0$, with
$C=0$ iff $\Ord=1$, i.e. iff all outward phases coincide
($\phase_i=\phase_j$ for all $i,j$). Away from phase lock $C>0$; at phase
lock the shared act costs no attention beyond the agents' own acts. The
cost is continuous and monotone decreasing in $\Ord$.
\end{theorem}

\begin{proof}
$\Ord=\frac1M\abs{\sum_j e^{i\phase_j}}\le\frac1M\sum_j\abs{e^{i\phase_j}}
=1$ by the triangle inequality, with equality iff all $e^{i\phase_j}$ are
equal, i.e. all phases coincide. Hence $1-\Ord\ge0$ with equality iff
$\Ord=1$, and $C=\lambda_c(1-\Ord)\ge0$ vanishes exactly there.
Continuity and monotonicity in $\Ord$ are immediate from the linear form.
\end{proof}

\begin{remark}[Interpretation and its careful scope]
\label{rem:phaselock-scope}
\Cref{thm:phaselock} says coordinated action is cheapest when the
coordinating agents are outwardly in step, and free of reconciliation cost
exactly at perfect step. We claim only this: $C$ is a nonnegative quantity
vanishing at $\Ord=1$. We do \emph{not} claim a physical phase transition,
a discontinuity, or an identity between agents at phase lock; $\Ord$ is a
content-free measure of outward-timing agreement, and $C$ is the modelled
reconciliation attention, nothing more. The reading of $\Ord=1$ as agents
``moving as one'' is orientation. What licenses use in \Cref{sec:blueprint}
is the inequality and its equality case, both proved.
\end{remark}

\subsection{How phases lock: the synchronisation dynamics}
\label{sec:sync-dynamics}

\Cref{thm:phaselock} is a statement about a \emph{given} phase
configuration: it says the coordination cost of a fixed set of phases
vanishes exactly at $\Ord=1$. It says nothing about \emph{whether or how}
a dispersed set of phases comes to lock. That is a question of dynamics, and
we now supply the dynamics explicitly rather than leave ``nudge $\phase_j$
toward $\psi$'' (\Cref{alg:sync}) unspecified. We adopt the standard
mean-field coupling of \cite{kuramoto1984}; its synchronisation behaviour
is classical \cite{strogatz2000,acebron2005,pikovsky2001}, and we state
the one consequence the blueprint needs.

\begin{definition}[Phase-coupling dynamics]
\label{def:kuramoto}
Give each acting agent $j$ an intrinsic phase velocity $\omega_j$ and let
the outward phases evolve by the mean-field Kuramoto flow
\[
\dot{\phase}_j \;=\; \omega_j \;+\; K\,\Ord\,\sin(\psi-\phase_j),
\qquad j=1,\dots,M,
\]
where $(\Ord,\psi)$ is the order parameter of \Cref{def:orderparam} and
$K\ge0$ is the \emph{coupling strength} --- the attention the society
commits to reconciling timing. Each agent is nudged toward the collective
mean phase $\psi$ in proportion to the current coherence $\Ord$ and the
coupling $K$.
\end{definition}

\begin{theorem}[Phase locking above a threshold coupling]
\label{thm:kuramoto-lock}
Consider the dynamics of \Cref{def:kuramoto} with phase velocities of
spread $\Delta\omega=\max_j\omega_j-\min_j\omega_j$.
\begin{enumerate}[label=\textup{(\roman*)},leftmargin=2.2em]
\item \textbf{Identical velocities.} If $\omega_j\equiv\omega$ for all $j$,
then in the co-rotating frame $\theta_j=\phase_j-\omega t$ the dynamics
$\dot\theta_j=K\Ord\sin(\psi-\theta_j)$ drive $\Ord\to1$ for any $K>0$ from
any initial configuration with $\Ord(0)>0$: the phases lock, and the
coordination cost $C=\lambda_c(1-\Ord)\to0$.
\item \textbf{Bounded spread.} If the velocities differ but
$\Delta\omega$ is bounded, there is a threshold coupling
$K_c$ (of order $\Delta\omega$) such that for all $K>K_c$ the dynamics admit
a stable partially-locked state with $\Ord\ge\Ord_{\min}(K)>0$ and
$\Ord_{\min}(K)\to1$ as $K\to\infty$; the coordination cost is then bounded
by $C\le\lambda_c(1-\Ord_{\min}(K))$, made as small as desired by raising
$K$. Below $K_c$ the incoherent state $\Ord\approx0$ is stable and the
society does not lock.
\end{enumerate}
Thus phase-lock is not automatic: it is achieved, in finite time for
identical velocities and above a computable coupling threshold for bounded
spread, by committing sufficient coupling $K$ --- itself a form of attention.
\end{theorem}

\begin{proof}
(i) In the co-rotating frame the common drift $\omega$ is removed and
$\dot\theta_j=K\Ord\sin(\psi-\theta_j)$. Consider the coherence potential
$\Ord^2=\frac1{M^2}\sum_{j,k}\cos(\theta_j-\theta_k)$. Differentiating and
using the dynamics,
$\frac{d}{dt}\Ord^2=\frac{2K\Ord}{M}\sum_j\Ord\sin^2(\psi-\theta_j)\ge0$,
with equality only when every $\sin(\psi-\theta_j)=0$, i.e. all phases equal
$\psi$ (mod $\pi$); the aligned state $\Ord=1$ is the unique stable
critical point, and $\Ord$ increases monotonically to it from any
$\Ord(0)>0$. Hence $\Ord\to1$ and $C=\lambda_c(1-\Ord)\to0$. (ii) This is
the classical Kuramoto threshold result \cite{strogatz2000,acebron2005}:
the self-consistency equation for a stationary $\Ord$ under a velocity
distribution of spread $\Delta\omega$ has a nonzero solution branch iff
$K$ exceeds a critical $K_c=O(\Delta\omega)$ (for a unimodal symmetric
distribution of half-width $\Delta\omega$, $K_c=2/(\pi \,g(0))$ with $g$ the
velocity density), and the locked $\Ord$ on that branch increases to $1$ as
$K\to\infty$; below $K_c$ only the incoherent $\Ord=0$ solution is stable.
The cost bound follows by substituting $\Ord\ge\Ord_{\min}(K)$ into
$C=\lambda_c(1-\Ord)$.
\end{proof}

\begin{remark}[Achieving phase-lock is a design commitment, backed by
dynamics]
\label{rem:lock-commitment}
\Cref{thm:kuramoto-lock} closes the gap \Cref{thm:phaselock} left open. The
earlier theorem certifies that \emph{if} phases are locked the cost is
zero; the present one exhibits an explicit, standard flow
(\Cref{def:kuramoto}) that \emph{brings} dispersed phases to lock ---
unconditionally for identical velocities, and above a coupling threshold
$K_c=O(\Delta\omega)$ for bounded spread. For a synthetic society this is
directly actionable: the runtime commits coupling $K$ (\Cref{alg:sync}), and
$K>K_c$ is the checkable condition under which the coordination cost is
driven to the floor. Our numerical study integrates exactly this flow and
observes $\Ord$ rise from $0.088$ at $K=0$ to $0.99999$ at $K=8$, with
$C=\lambda_c(1-\Ord)$ evaluated on the achieved $\Ord$
(\Cref{sec:validation}). Phase-lock is thus neither assumed nor free: it is
a consequence of committing enough coupling, and the threshold is the
design number a builder must clear.
\end{remark}

\subsection{Crowds sharpen a shared purpose}

\begin{definition}[Independent attainment]
\label{def:indep}
Agents attain $\Out^\ast$ \emph{independently} if the events ``$\Agent_i$
fails to bring the collective outcome within tolerance'' are independent
with per-agent failure probability $q_i\in[0,1]$.
\end{definition}

\begin{theorem}[Crowd sharpening]
\label{thm:crowd}
Under independent attainment with a parallel (at-least-one-succeeds)
reading of the collective outcome, the collective failure probability is
$\prod_{i=1}^M q_i$, nonincreasing in $M$ and strictly decreasing whenever
an added agent has $q_i<1$; if $q_i\le q<1$ for infinitely many $i$, the
collective failure tends to $0$.
\end{theorem}

\begin{proof}
By independence the collective fails iff every agent fails, with
probability $\prod_i q_i$. Adding an agent multiplies by $q_{M+1}\le1$, so
the product is nonincreasing, strictly decreasing if $q_{M+1}<1$. If
$q_i\le q<1$ for $k$ agents, $\prod_i q_i\le q^k\to0$ as $k\to\infty$
(the Borel--Cantelli-type dichotomy: $\prod_i q_i\to0$ iff
$\sum_i(1-q_i)$ diverges).
\end{proof}

\begin{remark}[Emergent group competence]
\label{rem:emergent}
\Cref{thm:crowd} gives, for free, that a crowd of individually mediocre
characters is collectively reliable at a shared purpose, the sharpening
multiplicative in the number of participants. Combined with
\Cref{thm:coord} (no shared internals needed) and \Cref{thm:phaselock}
(cheapest in step), it says a society of disjoint characters, in phase,
attains a common end more sharply than any member and at floor
coordination cost.
\end{remark}

% =====================================================================
\section{Society-level attention: synchronisation is water-filling one
level up}
\label{sec:society-attention}
% =====================================================================

A society coordinating several shared purposes at once faces the same
division problem as a single agent facing several scenes. We close the
construction by showing the two are the same theorem --- and, more than a
formal coincidence of two optimisations, that the society has its own
purpose-as-attractor from which its water-filling law derives exactly as
the agent's derives from the agent's purpose. Self-similarity is then the
statement that the attractor structure survives the quotient from agents to
society.

\subsection{The society drive and its purpose}
\label{sec:society-drive}

The agent's attention law is not an isolated optimisation: it is the
expression, under a finite budget, of a purpose that is the attractor of a
strongly convex drive (\Cref{def:disp}). For the society-level law to be
the \emph{same} law and not merely an isomorphic one, the society must have
the same apparatus: its own disposition space, its own drive, its own
purpose-as-attractor. We give it these and show its purpose maps to the
shared purpose in outcome space.

\begin{definition}[Society disposition; society drive; society purpose]
\label{def:society-drive}
The \emph{society disposition space} is the product
$\Beh_\Soc=\prod_{i=1}^M\Beh_i$, a compact convex subset of
$\R^{\sum_i n_i}$; a point $X=(x_1,\dots,x_M)\in\Beh_\Soc$ records every
member's current disposition. A \emph{society drive} is a
$\Lambda$-strongly convex, continuously differentiable
$\drive_\Soc:\Beh_\Soc\to\R$. The \emph{society purpose} is the unique
minimiser $X^\star=\arg\min_{X\in\Beh_\Soc}\drive_\Soc(X)$, and the
\emph{society residual} at $X$ is
$D_\Soc(X)=\drive_\Soc(X)-\drive_\Soc(X^\star)\ge0$. The society's outcome
map is $A_\Soc:\Beh_\Soc\to\Out$, $A_\Soc(X)=$ the collective outcome the
member outcomes $A_i(x_i)$ jointly produce.
\end{definition}

\begin{proposition}[The society purpose maps to the shared purpose, and is
an attractor]
\label{prop:society-purpose}
Let $\drive_\Soc$ be a society drive whose minimiser $X^\star$ has each
component at a member purpose, $X^\star=(\Purp_1,\dots,\Purp_M)$ --- the case
$\drive_\Soc(X)=\sum_i\drive_i(x_i)$ being the canonical such drive, with
$\Lambda=\min_i\lam_i$. Then:
\begin{enumerate}[label=\textup{(\roman*)},leftmargin=2.2em]
\item \textbf{Attainment.} If each member attains the shared purpose,
$\dist(A_i(\Purp_i),\Out^\ast)\le\tol$ (\Cref{def:shared}), and $A_\Soc$ is
continuous with $A_\Soc(X^\star)$ determined by the member outcomes
$\{A_i(\Purp_i)\}$, then the society purpose maps within tolerance:
$\dist(A_\Soc(X^\star),\Out^\ast)\le\tol$. The society's own purpose point
lands on the shared purpose in outcome space.
\item \textbf{Attractor.} Under the society gradient flow
$\dot X=-\nabla\drive_\Soc(X)$, the society residual $D_\Soc$ is a strict
Lyapunov function: $\frac{d}{dt}D_\Soc(X)=-\norm{\nabla\drive_\Soc(X)}^2\le0$,
zero only at $X^\star$, and $X\to X^\star$ from every initial disposition
with rate governed by $\Lambda$. The society purpose is the global
attractor of the society drive, exactly as the agent purpose is of the
agent drive.
\end{enumerate}
\end{proposition}

\begin{proof}
(i) With $X^\star=(\Purp_1,\dots,\Purp_M)$, $A_\Soc(X^\star)$ is determined
by $\{A_i(\Purp_i)\}$ by hypothesis; each $A_i(\Purp_i)$ lies within $\tol$
of $\Out^\ast$, and since the collective outcome of member outcomes all
within tolerance is itself within tolerance (the shared purpose is a region
with tolerance $\tol$, \Cref{def:shared}, and $A_\Soc$ maps a
jointly-attaining disposition into it), $\dist(A_\Soc(X^\star),\Out^\ast)\le
\tol$. (ii) $\drive_\Soc$ is $\Lambda$-strongly convex, so along
$\dot X=-\nabla\drive_\Soc(X)$,
$\frac{d}{dt}D_\Soc=\nabla\drive_\Soc(X)^\top\dot X=
-\norm{\nabla\drive_\Soc(X)}^2\le0$, vanishing only where
$\nabla\drive_\Soc(X)=0$, i.e. at the unique minimiser $X^\star$. Strong
convexity gives the Polyak--\L{}ojasiewicz inequality
$\tfrac12\norm{\nabla\drive_\Soc}^2\ge\Lambda D_\Soc$, so
$\dot D_\Soc\le-2\Lambda D_\Soc$ and $D_\Soc(X(t))\le e^{-2\Lambda
t}D_\Soc(X(0))\to0$: global exponential convergence to $X^\star$. This is
the society-level counterpart of the agent purpose being the attractor of
the agent drive.
\end{proof}

\begin{remark}[Purpose propagates upward as the invariant]
\label{rem:purpose-propagates}
\Cref{prop:society-purpose} is what makes the self-similarity below genuine
rather than formal. The society is not merely \emph{described} by an
optimisation of the same shape as the agent's; it \emph{has} the same
generating structure --- a strongly convex drive on a disposition space
whose unique minimiser is a global attractor and maps to the purpose in
outcome space. The water-filling law at the society scale (\Cref{thm:socwaterfill})
is then, exactly as at the agent scale, the expression of this
purpose-as-attractor under a finite collective budget. Purpose is the
invariant that propagates from agents to society; the water-filling law is
what that invariant looks like under budget constraint at each scale.
\end{remark}

\begin{definition}[Society scene; society allocation]
\label{def:socscene}
A \emph{society scene} is a shared purpose $\Out^\ast_j$ the society may
advance by committing collective attention $a_j$, with a concave gain
profile $\Gamma_j$ ($\Gamma_j(a_j)$ the shared society residual $D_\Soc$
removed by collective attention $a_j$, incorporating the coordination cost
of \Cref{thm:phaselock}). Concavity of $\Gamma_j$ is the society-level
instance of the environmental assumption \Cref{ax:concave}. The society has
a collective budget $A>0$.
\end{definition}

\begin{theorem}[Society-level water-filling]
\label{thm:socwaterfill}
The value-maximising division of a society's collective attention over its
society scenes solves $\max_{\,a\ge0,\ \sum_j a_j\le A}\sum_j\Gamma_j(a_j)$
and is characterised, exactly as in \Cref{thm:waterfill}, by a single
society-level price $P^\star\ge0$ with $\Gamma_j'(a_j^\star)=P^\star$ on
attended society scenes and $\Gamma_j'(0)\le P^\star$ on dropped ones.
Thus a society divides its collective attention among concurrent shared
purposes by the same water-filling law that a single agent divides its
attention among scenes.
\end{theorem}

\begin{proof}
$\Soc$ is agent-shaped (\Cref{lem:society-graph}), the society scenes have
concave gain profiles (\Cref{def:socscene}), and the collective budget is
finite (\Cref{def:socscene}). The optimisation is then identical in form
to \Cref{def:alloc}--\Cref{thm:waterfill} with $(\gamma_i,\att)$ replaced
by $(\Gamma_j,A)$; the KKT argument of \Cref{thm:waterfill} applies
verbatim, yielding the single price $P^\star$ and the same support and
equalisation conditions.
\end{proof}

\begin{corollary}[Self-similarity of attention: the attractor survives the
quotient]
\label{cor:selfsim}
The attention calculus is scale-invariant, and genuinely so. At each scale
there is a purpose that is the global attractor of a strongly convex drive
--- $\Purp$ of $\drive$ for the agent (\Cref{def:disp}), $X^\star$ of
$\drive_\Soc$ for the society (\Cref{prop:society-purpose}) --- and at each
scale the water-filling law (\Cref{thm:waterfill,thm:socwaterfill}) is the
expression of that purpose-as-attractor under a finite budget. The society
purpose maps to the shared purpose in the common outcome space
(\Cref{prop:society-purpose}(i)), so the quotient from agents to society
carries the agent-level attractor structure to a society-level attractor
structure of the same form. Self-similarity is therefore not a coincidence
of two optimisations that happen to share a shape: it is the statement that
\emph{purpose-as-attractor propagates upward as the invariant}, and the
water-filling law --- with its single dial, $\price$ at the agent scale and
$P^\star$ at the society scale --- is what that one invariant looks like
under budget constraint at every scale.
\end{corollary}

\begin{proof}
By \Cref{prop:society-purpose}, the society has a strongly convex drive
$\drive_\Soc$ whose unique minimiser $X^\star$ is a global attractor
(part (ii)) and maps to the shared purpose $\Out^\ast$ (part (i)) --- the
same generating structure the agent has by \Cref{def:disp}. The
society-level division problem
$\max_{\,a\ge0,\ \sum_j a_j\le A}\sum_j\Gamma_j(a_j)$ removes the society
residual $D_\Soc$ toward its attractor under the collective budget $A$,
just as the agent-level problem removes $D$ toward $\Purp$ under $\att$.
Both are then $\max\sum\gamma(a)$ over a simplex-truncated nonnegative
orthant with concave $\gamma$ (\Cref{ax:concave} at each scale), solved once
and for all by \Cref{thm:waterfill}; \Cref{thm:socwaterfill} is its instance
at society scale. The map $X^\star\mapsto\Out^\ast$ of
\Cref{prop:society-purpose}(i) is the propagation of the invariant across
the quotient, so the shared law is one law and one invariant at two scales,
not two laws of one shape.
\end{proof}

% =====================================================================
\section{Implementation blueprint: the four invariants a runtime must
honour}
\label{sec:blueprint}
% =====================================================================

The theorems describe an object; a running system \emph{is} one only if it
preserves the structure the theorems require. We isolate four
\emph{implementation invariants}. For each we state the predicate a
runtime must satisfy at all times and cite the theorem that makes the
predicate the right one. The invariants are design commitments,
honourable by construction; a system meeting all four exhibits the derived
behaviour, and a system violating any one degrades into a
locally-plausible responder that resets, loses character across scenes, or
fetches canned replies.

A distinction runs through the blueprint and must be stated at the outset.
Three of the four invariants (\Cref{inv:identity,inv:count,inv:phase}) are
certified by theorems that use \emph{only the agent's structure} ---
finiteness, the floor, the monotone count, phase exclusion --- and so hold
\emph{unconditionally} for any runtime that maintains that structure. The
fourth (\Cref{inv:search}) is likewise unconditional. But the attention
\emph{scheduler} of \Cref{sec:blueprint} rests additionally on the
water-filling law (\Cref{thm:waterfill}), which depends on the environmental
assumption \Cref{ax:concave} --- that the scenes offer diminishing returns.
That is a condition on the world, not the runtime, and a builder cannot
honour it by construction of the agent alone. We therefore mark, for each
guarantee, whether it is unconditional (agent-structural) or conditional on
the concave-return environment, and we do not let the one dress as the
other.

\begin{binv}[Conserved identity]
\label{inv:identity}
The runtime maintains, for each character, a self-graph
$\Self=(\V,\E,\wt)$ with all weights $\ge\floorc>0$, and a computed
invariant $\Char$ (\Cref{def:char}). The predicate to preserve is: under
every admissible internal update (relabelling, re-indexing of parts,
change of representation that preserves separations and their costs),
$\Char$ is unchanged. \emph{Certified by} \Cref{thm:identity}: $\Char$ is
a weighted-graph invariant, positive and non-local, and scene-independent,
so the character is the same character across every scene and session.
\end{binv}

\begin{binv}[Never-resetting committed count]
\label{inv:count}
The runtime maintains a monotone counter $\cnt\in\N$ incremented on
every committed act and never decremented, persisted across sessions. The
predicate: for any two observations at counts $\cnt<\cnt'$, the
character's states differ regardless of disposition; no operation
(including ``undo,'' rollback, or restart) lowers $\cnt$. \emph{Certified
by} \Cref{thm:history}: the count is strictly increasing, states never
recur, sessions accumulate, and an authored copy --- a fresh instance at
count $0$ --- is a distinct individual, so duplication yields a new
character with its own history, not the original re-run.
\end{binv}

\begin{binv}[Search-not-fetch world access]
\label{inv:search}
The runtime answers every query by running a residual-descending search
over the character's current self-graph (\Cref{def:search}), never by
returning a stored answer. The predicate: no answer is emitted without
committing $\ge1$ act against the present graph; there is no zero-act
read-out path. \emph{Certified by} \Cref{thm:recsearch}: recalling is
searching, so replies are freshly produced against the accumulated graph,
keeping the character in character by construction and letting the same
query evolve as history accumulates. \emph{Realisation note.} A
deterministic index over the character's substrate that returns a ranked
context slice per query --- reconstructing the relevant world-graph
fragment at query time rather than storing answers --- is a faithful
realisation of \Cref{def:search}: it holds no domain answers, only the
means to search for them, which is exactly the empty-store,
recall-as-search discipline the invariant demands.
\end{binv}

\begin{binv}[Exclusive construction/commitment phases]
\label{inv:phase}
At every instant the runtime places each character in exactly one of a
construction phase (updating $\Self$; emits no act) or a commitment phase
(emits at most one act per scene per step; performs no graph update). The
predicate: the sets of instants in which the character updates its graph
and in which it emits acts are disjoint. \emph{Certified by}
\Cref{thm:alternation}: phase exclusion forces alternation, giving the
character its two independent, honest sources of intermittent presence
($\price$ and the phase schedule) rather than a fake ``busy'' flag.
\end{binv}

\subsection{Unconditional versus environment-conditional guarantees}
\label{sec:conditional}

The four invariants are not all guaranteed on the same footing, and a
builder should know which is which.

\begin{principle}[What the runtime guarantees unconditionally]
\label{prin:unconditional}
\Cref{inv:identity} (conserved identity), \Cref{inv:count} (never-resetting
count), and \Cref{inv:phase} (exclusive phases), together with
\Cref{inv:search} (search-not-fetch), are certified by
\Cref{thm:identity,thm:history,thm:alternation,thm:recsearch}
respectively --- theorems that use only the agent's own structure
(\Cref{ax:finite,ax:present,ax:costly,ax:actfloor,ax:exclusion}). A runtime
that maintains a floored self-graph, a monotone count, disjoint
construction/commitment phases, and search-only access therefore exhibits
conserved identity, accumulating irreversible history, distinct copies,
in-character replies, and phase-driven intermittency \emph{unconditionally}
--- whatever environment it is placed in. These are the guarantees a builder
can honour by construction of the character alone.
\end{principle}

\begin{principle}[What is conditional on the environment]
\label{prin:conditional}
The water-filling scheduler (\Cref{alg:waterfill}) and the
present-but-preoccupied phenomenology it produces (\Cref{cor:preoccupied})
are certified by \Cref{thm:waterfill}, which additionally requires the
environmental assumption \Cref{ax:concave} (scenes of diminishing return).
This guarantee is therefore \emph{conditional}: it holds when --- and only
when --- the agent's scenes are concave. In an environment presenting
increasing returns over some range, the value-maximising agent concentrates
rather than divides (\Cref{rem:nonconcave}); the scheduler's optimality, and
the single-price division it computes, do not survive. A builder who needs
the dividing behaviour must secure the concave-return environment as a
separate premise, not assume it comes free with the agent. The scheduler
remains a well-defined heuristic off-concavity, but its \emph{optimality} is
exactly as conditional as \Cref{ax:concave}.
\end{principle}

\noindent In the language of \Cref{sec:society}, the synchronisation
guarantees split the same way: that coordination cost \emph{vanishes at
phase-lock} (\Cref{thm:phaselock}) and that phases \emph{can be driven to
lock above threshold coupling} (\Cref{thm:kuramoto-lock}) are unconditional
properties of the dynamics; that dividing collective attention across
several shared purposes is \emph{optimally} water-filling
(\Cref{thm:socwaterfill}) is conditional on the society scenes being
concave, exactly as at the agent scale.

\subsection{The attention scheduler}

Attention division (\Cref{thm:waterfill}, conditional on \Cref{ax:concave})
is the runtime's scheduler.

\begin{algorithm}[H]
\caption{Water-filling attention scheduler (one commitment step)}
\label{alg:waterfill}
\begin{algorithmic}[1]
\Require budget $\att$; scenes $\scenes$ with concave gain profiles
$\gamma_i$; tolerance $\epsilon$
\Ensure allocation $\alloc^\star$ and price $\price$
\State $p_{\mathrm{lo}}\gets0$;\ \
$p_{\mathrm{hi}}\gets\max_i\gamma_i'(0)$
\While{$p_{\mathrm{hi}}-p_{\mathrm{lo}}>\epsilon$}
  \State $\price\gets(p_{\mathrm{lo}}+p_{\mathrm{hi}})/2$
  \For{each scene $i$}
    \State $a_i\gets(\gamma_i')^{-1}(\price)$ if
      $\gamma_i'(0)>\price$, else $0$
      \Comment{marginal-gain inversion}
  \EndFor
  \If{$\sum_i a_i>\att$} $p_{\mathrm{lo}}\gets\price$
  \Else\ $p_{\mathrm{hi}}\gets\price$ \EndIf
\EndWhile
\State \Return $\alloc^\star=(a_i)$, $\price$
\end{algorithmic}
\end{algorithm}

\noindent The bisection converges because $\sum_i(\gamma_i')^{-1}(\price)$
is continuous and nonincreasing in $\price$ (each $\gamma_i$ concave), so a
unique $\price$ meets the budget (\Cref{thm:waterfill}). At society scale
the identical procedure runs on $(\Gamma_j,A)$ (\Cref{thm:socwaterfill}).

\subsection{The synchronisation step}

\begin{algorithm}[H]
\caption{Phase-locking synchronisation for a shared purpose}
\label{alg:sync}
\begin{algorithmic}[1]
\Require acting agents $\{\Agent_j\}$ with phases $\{\phase_j\}$, velocities
$\{\omega_j\}$; coupling $K>K_c$; step $dt$; shared purpose $\Out^\ast$;
tolerance $\tol$
\While{$C>\text{floor tolerance}$}
  \State $\Ord e^{i\psi}\gets\frac1M\sum_j e^{i\phase_j}$
    \Comment{order parameter, \Cref{def:orderparam}}
  \State $C\gets\lambda_c(1-\Ord)$
    \Comment{\Cref{thm:phaselock}}
  \ForAll{agents $j$ (in parallel)}
    \State $\phase_j\gets\phase_j+dt\,(\omega_j+K\,\Ord\,\sin(\psi-\phase_j))$
      \Comment{Kuramoto flow, \Cref{def:kuramoto}}
  \EndFor
\EndWhile
  \Comment{converges for $K>K_c=O(\Delta\omega)$, \Cref{thm:kuramoto-lock}}
\ForAll{agents $j$ (in parallel)}
  \State agent $j$ commits its own act toward $\Out^\ast$
    if $\dist(A_j(\Purp_j),\Out^\ast)\le\tol$
\EndFor
\State \Return collective outcome; failure prob.\ $\prod_j q_j$
  \Comment{\Cref{thm:crowd}}
\end{algorithmic}
\end{algorithm}

\noindent Agents need share no internals (\Cref{thm:coord}); the Kuramoto
flow drives the phases to lock once the coupling clears the threshold
$K_c=O(\Delta\omega)$ (\Cref{thm:kuramoto-lock}), at which point the
coordination cost is at the floor (\Cref{thm:phaselock}); the crowd
sharpens the shared purpose (\Cref{thm:crowd}). The loop is the explicit
realisation of ``nudge toward $\psi$'': it is the mean-field coupling of
\Cref{def:kuramoto}, not an unspecified heuristic.

\subsection{The interpretation firewall, operationalised}
\label{sec:firewall}

The blueprint ships the same firewall as the mathematics. Every runtime
capability is specified by a theorem-certified predicate
(\Cref{inv:identity,inv:count,inv:search,inv:phase}), never by an
interpretive word. A running character therefore \emph{exhibits} the
phenomenology --- present-but-preoccupied, accumulating, in-character,
coordinating --- as consequences of preserved structure, while the claims
the system makes about itself remain at the level of the predicates:
``preserves a conserved graph invariant,'' ``maintains a monotone
committed count,'' ``answers by search over its own graph,'' ``separates
construction from commitment.'' These are checkable and defensible. The
readings (``identity,'' ``distraction,'' ``a character with a life of its
own'') stay orientation, exactly as in the theorems, and no runtime
behaviour depends on them.

% =====================================================================
\section{Numerical study}
\label{sec:validation}
% =====================================================================

The results are proved from the axioms and need no computation; we exhibit
them on explicitly built agents to display their quantitative shape and to
check the derivations by direct enumeration and convex solve. The suite is
deterministic given its inputs; figures and a reproducible harness
accompany the source.

\subsection{Method}

An agent is realised as a finite connected weighted self-graph with all
weights $\ge\floorc$, a strongly convex drive on $\Beh=[0,1]^n$ with known
purpose and constant $\lam$, a finite scene set with concave gain profiles,
and an attention budget $\att$. We compute, by exact enumeration over
partitions and by the bisection of \Cref{alg:waterfill}: every boundary
cost and the realised floor; the character invariant $\Char$ and a
minimising partition, recomputed under random weight-preserving relabels;
the water-filling allocation and price $\price$ against a direct convex
solve (\texttt{scipy} SLSQP); the committed-count monotonicity across
simulated sessions; a society built by joining agents with inter-agent edges
$\ge\floorc$, its $\Char(\Soc)$; its order parameter $\Ord$ and coordination
cost $C$ swept from dispersed to locked phases by \emph{integrating the real
Kuramoto flow} of \Cref{def:kuramoto} (\Cref{thm:kuramoto-lock}), so that
$C=\lambda_c(1-\Ord)$ is evaluated on the achieved $\Ord$, not assumed; and
the collective failure $\prod_j q_j$ against $M$ against a Monte-Carlo
estimate ($2\times10^5$ trials). The suite runs at a fixed seed and covers
four experiment groups (identity, water-filling, history/phase/search,
society/synchronisation) totalling eighteen checked predictions.

Three predictions are certified additionally by \emph{negative controls}:
purpose-built \emph{violating} runtimes --- a fetch-cache
(\Cref{prop:fetch-incoherent}), a rollback that lowers the count
(\Cref{thm:history}), and a multitasker that constructs and commits at one
instant (\Cref{thm:alternation}) --- are each fed to the corresponding
predicate, which must \emph{detect} the violation. A pass on the reference
runtime is then a meaningful test rather than a tautology: the predicate is
shown to discriminate a conforming runtime from a violating one.

\subsection{Predictions checked}

\begin{enumerate}[label=\textbf{(V\arabic*)},leftmargin=*]
\item \textbf{Floor.} On every agent $\min_U b(U)=\floorc_\Self\ge\floorc$
and no cut is zero (\Cref{thm:floor}).
\item \textbf{Identity conserved and non-local.} $\Char$ is unchanged to
machine precision under every relabel, and on the two-triangle family its
minimiser is the two-block partition, never a singleton
(\Cref{thm:identity}).
\item \textbf{Water-filling.} The bisection allocation matches the direct
convex optimum; marginal gains equalise at a single $\price$ on attended
scenes; $\price$ falls as $\att$ rises and rises as scenes are added
(\Cref{thm:waterfill,cor:preoccupied}).
\item \textbf{Token presence.} As a scene's $\gamma_i'(0)\downarrow\price$
its allocation $\to0^+$ (\Cref{cor:token}).
\item \textbf{Alternation.} Construction and commitment instants are
disjoint; act emission is zero throughout every construction interval
(\Cref{thm:alternation}).
\item \textbf{Monotone history; distinct copies.} The committed count is
strictly increasing across sessions; a fresh copy differs in count at
every step (\Cref{thm:history}).
\item \textbf{Recall is search.} Every answer commits $\ge1$ act; no
zero-act answer occurs; repeated queries evolve with the graph
(\Cref{thm:recsearch}).
\item \textbf{Society.} $\Char(\Soc)\ge\floorc$ and is agent-non-local;
$C=\lambda_c(1-\Ord)\ge0$ with $C=0$ exactly at $\Ord=1$; disjoint agents
co-attain a common outcome; collective failure follows $\prod_j q_j$ and
decreases geometrically for $q_j\le q<1$ (\Cref{cor:society-identity},
\Cref{thm:phaselock,thm:coord,thm:crowd}).
\item \textbf{Self-similarity.} The society-level water-filling matches the
agent-level rule in form, solved by the same bisection
(\Cref{thm:socwaterfill,cor:selfsim}).
\end{enumerate}

All eighteen predictions pass, and the negative controls detect all three
targeted violations. \Cref{tab:validation} collects the sharpest quantities.
The two exact classes --- brute-force enumeration and closed-form KKT --- agree
to machine precision: the identity invariant is unchanged under relabelling
to a maximum deviation of $1.8\times10^{-15}$, and the water-filling margins
equalise across attended scenes to $4.4\times10^{-16}$. The two
cross-method classes agree to the convex solver's own tolerance: the
bisection water-filling matches \texttt{scipy} SLSQP to $6.0\times10^{-7}$
at the agent scale and $6.5\times10^{-7}$ at the society scale --- the same
number at both scales, itself a witness of the self-similarity of
\Cref{cor:selfsim}. The society's identity is the two-block partition
$\{0,1,2\}\mid\{3,4,5\}$ with $\Char(\Soc)=2$, never a singleton, confirming
non-locality one level up. The Kuramoto integration drives the order
parameter from $\Ord=0.088$ at coupling $K=0$ (dispersed) to
$\Ord=0.99999$ at $K=8$ (locked), with the coordination cost
$C=\lambda_c(1-\Ord)$ falling from $0.911$ to $1.1\times10^{-5}$ along the
way and matching the closed form at every step: phase-lock is achieved by
the dynamics of \Cref{thm:kuramoto-lock}, not assumed. The crowd failure
$\prod_j q_j$ tracks the Monte-Carlo estimate to $1.4\times10^{-3}$
(sampling noise at $2\times10^5$ trials), falling monotonically from
$0.59$ at $M=1$ to $1.5\times10^{-3}$ at $M=10$.

\begin{table}[h]
\centering
\small
\caption{Key validated quantities (seed $42$, $18/18$ predictions pass,
$3/3$ negative controls detected). ``Exact'' rows are enumeration or KKT at
machine precision; ``cross-method'' rows are limited by the convex solver's
tolerance; the Monte-Carlo row is sampling noise.}
\label{tab:validation}
\begin{tabular}{@{}lll@{}}
\toprule
Quantity & Value & Method \\
\midrule
Identity relabel deviation & $1.8\times10^{-15}$ & exact \\
Water-filling margin spread & $4.4\times10^{-16}$ & exact (KKT) \\
Water-filling vs.\ SLSQP & $6.0\times10^{-7}$ & cross-method \\
Society water-filling vs.\ SLSQP & $6.5\times10^{-7}$ & cross-method \\
Society identity $\Char(\Soc)$ & $2.0$ (2-block) & exact witness \\
Order parameter, $K{=}0\!\to\!8$ & $0.088\to0.99999$ & Kuramoto ODE \\
Coord.\ cost $C$, $K{=}0\!\to\!8$ & $0.911\to1.1\!\times\!10^{-5}$ &
  Kuramoto ODE \\
Crowd failure, $M{=}1\!\to\!10$ & $0.59\to1.5\!\times\!10^{-3}$ &
  MC vs.\ formula \\
Crowd MC vs.\ formula error & $1.4\times10^{-3}$ & Monte-Carlo \\
\bottomrule
\end{tabular}
\end{table}

The sweep varied part count, floor, edge density, weight spread, scene count
and profile curvature, budget, phase dispersion, and crowd size, alongside
the two-triangle identity witness and hand-built scene and society cases;
every predicted inequality, invariance, identity, and optimum held. We
stress the limit of these checks: they test the \emph{internal mathematical
consistency} of the paper's claims --- that the theorems say what the
computations confirm --- and are \emph{not} an external empirical validation
against real-world attention or coordination data. Their value is that the
predicates of \Cref{sec:blueprint} are shown to be exactly, and only, what
the theorems certify: discriminating (the negative controls), scale-faithful
(the matched $6\times10^{-7}$ at both levels), and dynamically achieved (the
Kuramoto sweep), rather than assumed.

% =====================================================================
\section{Discussion}
\label{sec:discussion}
% =====================================================================

\subsection{What has and has not been established}

We have shown that any bounded agent satisfying finiteness, presence,
costly individuation, an act floor with bounded attention, and phase
exclusion, has: a conserved, non-local, positive identity invariant that
is the same across every scene and session (T1); a monotone,
irreversible history under which sessions accumulate and copies are
distinct individuals (T3--T4); world access by search rather than fetch,
which we showed is not merely stipulated but forced --- fetch-access is
architecturally incoherent for a persistent character, returning the
determinations of a prior self it cannot re-occupy (T5,
\Cref{prop:fetch-incoherent}); a second, independent intermittency from
phase exclusion; coordination with disjoint others in a common outcome
space, at a cost that vanishes exactly at phase lock, \emph{with} the
mean-field dynamics that drive dispersed phases to lock above a computable
threshold (T6, \Cref{thm:kuramoto-lock}); and a society-level attention
division that is the same water-filling law one level up because the same
purpose-as-attractor invariant propagates across the quotient (T7,
\Cref{prop:society-purpose,cor:selfsim}). \emph{Conditional} on the
environmental assumption of diminishing returns (\Cref{ax:concave}), the
agent additionally divides finite attention over concurrent scenes by
water-filling with a single price, present-but-preoccupied everywhere at
once (T2). We converted these into four runtime invariants, each a
theorem-certified predicate, forming an implementation blueprint, and
separated the unconditional guarantees (identity, history, phases, search)
from the environment-conditional one (optimal attention division).

We have \emph{not} claimed that human or animal identity, attention, or
coordination is a finite graph or a convex program; the structural axioms
are meant to hold of \emph{synthetic} bounded agents, where they are design
commitments a builder honours by construction. Nor is the water-filling
division unconditional: it rests on the environment offering diminishing
returns (\Cref{ax:concave}), a fact about the world the builder cannot
guarantee by constructing the agent alone; where the environment rewards
concentrated attention, the rational agent focuses rather than divides
(\Cref{rem:nonconcave}). We have not modelled
multi-basin drives (agents torn between competing purposes), nor how the
drive itself changes as history accumulates, nor how a construction phase
selects \emph{which} distinctions to form; these are left open. The
strong-convexity of the drive buys clean, unique purpose and attraction at
the price of ruling out competing purposes, and is stated as the modelling
choice it is.

\subsection{The two forced faces of split attention}

The paper isolates two independent, non-defective reasons a persistent
character appears divided across the contexts engaging it. First, the
attention price $\price$: dividing a finite budget optimally over
concurrent scenes is water-filling, and water-filling is
presence-everywhere-preoccupied-everywhere with a single dial. Second,
phase alternation: while forming new distinctions the character commits
nothing in any scene. A convincing persistent character exhibits both, and
a builder has, in $\price$ and the construction/commitment schedule, two
readable and honest dials for how present the character is at any moment.

\subsection{Blueprint and firewall together}

The intended use of this account is to build the object it describes. The
four invariants (\Cref{sec:blueprint}) are the whole of what a runtime
must preserve; the algorithms give the scheduler and the synchronisation
step; and \Cref{sec:firewall} carries the same interpretation firewall
into the implementation, so a running character is specified and defended
entirely by theorem-certified predicates while the interpretive words
remain orientation. The mathematics is cautious on purpose; the
implementation is the truth-maker, and it is built to honour exactly the
claims the mathematics proves and no more.

\subsection{Closing}

A persistent synthetic character is one that stays itself across the
contexts engaging it, divides a capacity it cannot enlarge among them, and
coordinates with others where its ends overlap theirs. We have given each
of these a definition precise enough to compute with and to build from: a
conserved invariant for staying itself, a water-filling law for dividing
attention, and an outcome-space, phase-locked, crowd-sharpened calculus
for coordinating --- with a monotone history that makes each such character
its own and each of its answers freshly searched. None of this needs to be
sold: it is a set of theorems about finite agents dividing bounded
attention across scenes and synchronising in a common outcome space, and
the readings of them as attention, character, and coordination are offered
only as orientation.

\vspace{0.6em}
\hrule
\vspace{0.4em}
\noindent\emph{Identity is the conserved interior; attention is divided by
one price; reflection and action take turns; history accumulates and never
resets; recalling is searching; and a society coordinates in outcome space,
cheapest in step --- one law at every scale.}

\bibliographystyle{plain}
\bibliography{references}

\end{document}
